\section{A linear rank cutoff over proper prime powers}
\label{sec:proper-prime-power-linear-cutoff}

\paragraph{Status and provenance.}
D-TD66, question Q2fv, is an accepted extension of D-TD65.
The complete scalar-field genus argument below is independent of
the prime-field exception table in Frohardt--Guralnick--Magaard.
Its spherical translation strategy and the actual source interfaces
are shared accepted inputs. The weak scalar-choice Scott inequality
is used over the actual field, including fields of orders
\(4,8,9,16\). No binary scalar-trivial strengthening is imported.
The author has read the complete proof, invoked source interfaces,
both exact supplements, all outputs and their retained histories.
Complete root and independent whole-proof reviews pass;
final documentary bindings and histories also pass.

Put
\[
 D_0=10^{64},\quad T_0=2^{1100000},\quad B_0=10^{30},
 \quad C_0=2\cdot10^{30},\quad J_0=2C_0+1.
 \tag{QPL.1}
\]
A projective support always allows every target eigenvalue.

\begin{theorem}[D-TD66]
Let an actual connected finite \(J\)-Galois sphere cover have
faithful irreducible complex coefficient module \(V\), of degree
\(e\), and genuine generating product-one inertia with
\[
                  \sum_i\operatorname{rank}(1-T_i)=2e+6.
 \tag{QPL.2}
\]
Suppose an actual stable whole factorization \(V=W\otimes U\)
has a finite projective image on \(W\) containing and normalizing
a faithful irreducibly acting perfect central product
\(E_1\cdots E_t\). Its quasisimple components are permuted
transitively. Each is an actual central quotient of
\(\operatorname{SL}_d(q)\), where \(q=p^f\), \(p\) is prime,
\(f\ge2\), and \(d\ge D_0\). Write
\(W=\bigotimes_{j=1}^tW_j\), with common component degree \(a\).
If \(t\le T_0\), this configuration is impossible.

The entire partner \(U\), including its image, dimension and
scalar phases, is unrestricted. There is no bound on \(q\) or
\(f\), no projected-mark minimum, and no assumption that a
component-stabilizer base is rational.

Consequently, in exactly the primitive minimal naturally
nonsplit exact-imaginary-quadratic source class of D-TD54,
every layer of a fixed complete D-TD12 recursion having this
named-cover presentation has \(d<D_0\). Combined with D-TD65,
the same rank cutoff holds for all prime powers in this
linear named-cover scope.
\end{theorem}

\subsection*{1. A scalar-field genus estimate}

We prove a slightly more general elementary input. Let \(Q\ge4\)
be a prime power and \(d\ge1000\). Let
\(G\subset\operatorname{PGL}_d(Q)\) contain a normal nonabelian
simple subgroup \(S\), of order greater than 60, whose linear
preimage acts irreducibly on the natural space.
For \(x\in G\), let \(v(x)\) be \(d\) minus the largest rational
\(\mathbf F_Q\)-eigenspace dimension of a scalar lift, using
dimension zero when no rational eigenvalue occurs.
This is unchanged by multiplying the lift by an \(\mathbf F_Q\)
scalar. Let a nonempty finite \(G\)-set \(\Omega\), possibly
with several orbits, have degree \(N\) and satisfy
\[
       f_\Omega(x)\le Q^{-v(x)}+\frac1{10000}
                     \qquad(x\ne1).
 \tag{QPL.3}
\]
For an actual connected \(G\)-cover of a sphere, let \(Y_\alpha\)
be the point-stabilizer quotient curves for the orbits of \(\Omega\).
We claim that \(\sum_\alpha(g(Y_\alpha)-1)\ge N/10000\).
The FPR hypothesis makes the \(G\)-action on \(\Omega\) faithful:
a nonidentity element in its kernel would have fixed fraction one,
whereas the right side of (QPL.3) is at most \(1/4+1/10000\).
In the transitive case the conclusion is \(g(Y)\ge1+N/10000\).

Remove identity entries and write their exact projective orders
as \(o_1,\ldots,o_r\). The regular \(G\)-cover has genus at least
two by the complete genus-zero/elliptic argument of D-TD02,
so \(A=\sum_i(1-1/o_i)>2\) and \(r\ge3\).
For \(x\) of order \(o\), set
\(\kappa(x)=o^{-1}\sum_{j=0}^{o-1}Q^{-v(x^j)}\).
A nonidentity projective power has \(v\ge1\). Thus
\[
 \kappa(x)\le\frac1o+\frac{o-1}{Qo}.
 \tag{QPL.4}
\]
Suppose the claimed summed-genus bound fails. The possibly
disconnected Riemann--Hurwitz equality is
\(\sum_i\operatorname{ind}_\Omega(x_i)
=2N+2\sum_\alpha(g(Y_\alpha)-1)\).
Cauchy--Frobenius averaging and (QPL.3) therefore imply
\[
 \begin{gathered}
 \sum_i\kappa(x_i)>r-2-\frac{A}{10000}-\frac1{5000},\\
 (3/4-1/10000)A<2+1/5000,\qquad
 A<27/10,\qquad 3\le r\le5,\\
                  \sum_i\kappa(x_i)>r-2-\frac1{1000}.
 \end{gathered}
 \tag{QPL.5}
\]
These estimates use \(Q\ge4\), not its characteristic.

There is also a useful order bound with no minimal-degree table.
If \(v(y)=1\), rescale a lift to \(I+uv^T\).
Its linear order equals its projective order because its
nonzero fixed eigenspace makes any scalar power equal to one.
If \(v^Tu=0\), that order is the characteristic \(p\).
Otherwise the matrix is semisimple with its other eigenvalue
in \(\mathbf F_Q^*\), so the order divides \(Q-1\).
In one cyclic group at most \((p-1)+(Q-2)\le2Q-3\)
nonidentity powers can therefore have support one.
All the other nonidentity powers have support at least two.
Consequently
\[
                         \kappa(x)<\frac3{o(x)}+\frac1{Q^2}
                                  \le\frac3{o(x)}+\frac1{16}.
 \tag{QPL.6}
\]
This includes every proper prime power, without treating
\(Q\) as a prime.

\paragraph{The weak system inequality.}
For a spherical signature \(\mathbf e\), either all entries
one or a permutation of
\((m,m,1,\ldots)\), \((2,2,m,1,\ldots)\), or
\((2,3,m,1,\ldots)\) with \(m=3,4,5\) in the last case, put
\(L=2/(2-\sum_i(1-1/e_i))\) and \(C_i=L/e_i\).
The universal spherical cover gives a product-one tuple of
\(C_i\) conjugates of \(x_i^{e_i}\).
Its generated normal subgroup has spherical \emph{quotient}
of \(G\), not necessarily the universal group's exact order.
That normal subgroup contains \(S\): a simple group of order
greater than 60 cannot map nontrivially into these solvable
or order-at-most-60 quotients.

Use the direct Scott proof supplied in D-TD65: for a
product-one linear tuple \(g_j\), the map
\(\bigoplus\operatorname{im}(g_j-1)\to\mathbf F_Q^d\),
\((w_j)\mapsto\sum_jg_1\cdots g_{j-1}w_j\), has image
codimension equal to the dual invariant dimension and its
kernel contains the image of \(v\mapsto((g_j-1)v)_j\).
The generated projective group contains the irreducible
\(S\), so both invariant dimensions are zero.

Choose all but one lifted generator to attain their largest
rational eigenspace; use the last scalar to make their
product one. The last rank is at most \(d\). For any selected
index \(i_*\), this proves
\[
 (C_{i_*}-1)v(x_{i_*}^{e_{i_*}})
             +\sum_{i\ne i_*}C_i v(x_i^{e_i})\ge d.
 \tag{QPL.7}
\]
All scalar choices here are in \(\mathbf F_Q^*\).
Only this weak inequality is used, even in characteristic two.

We record its exact consequences. Taking all \(e_i=1\) gives
\(\sum_{i\ne j}v(x_i)\ge d\) for every \(j\).
For three entries, the signature \((o_i,o_i,1)\),
suitably permuted, gives \(o_i v(x_j)\ge d\) when \(i\ne j\).
For four entries, if
\(v(x_1)+v(x_2)\le d/8\), applying
\((1,1,o_3,o_3)\) and then
\((1,1,o_4,o_4)\) shows \(o_3,o_4\ge8\).
For a signature \((2,u,v)\), use \((2,2,v)\)
and \((2,u,2)\) to obtain
\[
 v(x_2^2)\ge\frac d{v-1},\qquad
 v(x_3^2)\ge\frac d{u-1}.
 \tag{QPL.8}
\]
For \((2,3,o)\), applying \((2,3,k)\),
\(k=2,3,4,5\), with \(i_*=3\), gives respectively
\(v(x_3^k)\ge d/2,d/3,d/5,d/11\).

If \(\gcd(j,o)=1\), a lift of \(x\) is a scalar times a
power of a lift of \(x^j\), and conversely. Their rational
eigenspace dimensions are equal. More generally powers with
the same gcd of exponent and order have the same \(v\).
Put
\[
 \delta=Q^{-d/29}<10^{-6},\qquad
 g(o)=\frac1o+\frac{o-1-\varphi(o)}{4o}.
 \tag{QPL.9}
\]
If \(v(x)\ge d/29\), its generator powers have the same
support and
\(\kappa(x)\le g(o)+\delta\).
Here \(g(2)=1/2\), \(g(o)\le1/3\) for \(o\ge3\), and
\(g(o)\le5/16\) for \(o\ge4\).
These last two bounds follow from \(\varphi(o)\ge2\)
for \(o\ge3\).

\paragraph{Five entries.}
The bound \(A<27/10\) forces every order to be two or three:
four orders two and one order at least four already give
\(A\ge11/4\).
The system inequality with each index omitted makes the
two largest supports at least \(d/4\). For orders two or
three, their \(\kappa\)'s are at most \(1/2+\delta\).
The other three are at most \(5/8\) by (QPL.4). Hence
\[
                    \sum_i\kappa(x_i)\le23/8+2\delta<2.999,
 \tag{QPL.10}
\]
contradicting (QPL.5).

\paragraph{Four entries.}
Sort the supports increasingly. If \(v_1+v_2\le d/8\),
the last two orders are at least eight, so (QPL.4) gives
\(\sum_i\kappa_i\le2(5/8)+2(1/4+3/32)=31/16<1.999\).
Otherwise the last three supports exceed \(d/16\).
If one of these entries is not an involution, their bounds
give \(\sum_i\kappa_i\le47/24+3\delta<1.999\).
If all three are involutions, the remaining order is at
least three because \(A>2\). The spherical signature
\((1,2,2,2)\), with the first index selected, forces the
remaining support to be at least \(d/3\). Then
\(\sum_i\kappa_i\le11/6+4\delta<1.999\).
Every case contradicts (QPL.5).

\paragraph{Three entries.}
Sort the orders \(o_1\le o_2\le o_3\).
Equations (QPL.4)--(QPL.6) give
\[
 .999<\sum_i\kappa_i<\frac34+\frac6{o_2},
               \qquad o_2\le24.
 \tag{QPL.11}
\]
All three generator supports are therefore at least \(d/24\).
If \(o_1\ge3\), the signature is not \((3,3,3)\), since
\(A>2\). Thus
\(\sum_i\kappa_i\le47/48+3\delta<.999\), a contradiction.

We may assume \(o_1=2\), so \(o_2\ge3\).
If \(o_3\ge30\), the bounds on the first two generators and
(QPL.6) give
\(\sum_i\kappa_i<239/240+2\delta<.999\).
Thus \(o_3\le29\).
By (QPL.8), the squares of \(x_2,x_3\) also have support
at least \(d/29\). Their orders exceed two, so these squares
are not projectively trivial.
It follows that
\[
 \begin{gathered}
 \kappa(x_j)\le h(o_j)+\delta\quad(j=2,3),\\
 h(o)=\frac{1+(o-1-\varphi(o)
                   -{\boldsymbol 1}_{2\mid o}\varphi(o/2))/4}{o},\\
 h(o)\le1/4\ (4\le o\le29),\qquad
 h(o)\le5/24\ (5\le o\le29).
 \end{gathered}
 \tag{QPL.12}
\]
These are finite integer inequalities. To make the complete
calculation explicit, the values of
\(\varphi(o)+{\boldsymbol 1}_{2\mid o}\varphi(o/2)\)
for \(o=5,\ldots,29\), in order, are
\[
 \begin{split}
 &4,4,6,6,6,8,10,6,12,12,8,12,16,12,18,\\
 &12,12,20,22,12,20,24,18,18,28.
 \end{split}
 \tag{QPL.13}
\]
Each satisfies \(6(\cdot)-o-18\ge0\), with minimum zero at
\(o=6\), proving the \(5/24\) bound. Directly \(h(4)=1/4\).

If \(o_2\ge4\), positivity excludes \((2,4,4)\), so
(QPL.12) gives
\(\sum_i\kappa_i\le23/24+3\delta<.999\).

Finally let \(o_2=3\); then \(7\le o_3\le29\).
The four spherical signatures following (QPL.8), together
with the generator support, control every power of \(x_3\)
whose exponent has gcd \(1,2,3,4,\) or \(5\) with \(o_3\).
All these powers have support at least \(d/29\).
For \(14\le o_3\le29\), any remaining power has order at
most four. There are at most
\(\varphi(2)+\varphi(3)+\varphi(4)=5\) such nonidentity
powers. Thus \(\kappa(x_3)\le9/(4o_3)+\delta\le9/56+\delta\).
For \(o_3=7,8,9,10,11,13\) every nonidentity power is
already controlled. For \(o_3=12\) only the order-two
power may remain, giving \(5/48+\delta\).
All seven values satisfy the same \(9/56+\delta\) bound.
Consequently
\[
                \sum_i\kappa_i\le
                   \frac12+\frac13+\frac9{56}+3\delta
                    =\frac{167}{168}+3\delta<.999.
 \tag{QPL.14}
\]
This final contradiction proves the claimed general sphere lemma.

The same summed-genus bound holds over every base genus.
For an elliptic base there must be nontrivial branching, since
an unramified finite elliptic monodromy group is abelian.
The action is faithful, and (QPL.3) bounds a nonidentity
branch fixed fraction by \(1/4+1/10000<2/3\). Its index is
at least \(N/6\), since every moved permutation orbit has
length at least two. Riemann--Hurwitz then gives
\(\sum_\alpha(g(Y_\alpha)-1)\ge N/12\).
For base genus at least two it gives this sum at least \(N\).
These statements use the actual connected regular \(G\)-cover,
but neither require \(\Omega\) to be transitive nor assume
that an intermediate base is rational.

\subsection*{2. Projective points and every intermediate base}

Apply the lemma to
\(\operatorname{PSL}_d(q)\le G\le\operatorname{PGL}_d(q)\),
with \(q\ge4\), \(d\ge1000\), and its projective points.
The scalar preimage of PSL acts irreducibly.
If the largest rational eigenspace has dimension \(k\ge d/2\),
all other rational eigenspaces have total dimension at most
\(d-k\). Convexity of \(q^u\) therefore gives
\[
 f_{\rm pt}(x)\le
      \frac{q^k+q^{d-k}-2}{q^d-1}
      \le q^{k-d}+2q^{-d/2}.
 \tag{QPL.15}
\]
If \(k<d/2\), there are at most \(q-1\) rational eigenvalues;
counting their lines gives the weaker but sufficient bound
\(f_{\rm pt}(x)\le2q^{1-d/2}\).
Thus always \(f_{\rm pt}(x)\le q^{-v(x)}+2q^{1-d/2}\),
and \(2q^{1-d/2}<1/10000\).
The general lemma applies. This direct proof does not transfer
any prime-field numerical exception table.

For a connected point cover \(Y\to X\) of degree
\(N=(q^d-1)/(q-1)\) and the indicated linear monodromy, we have
\[
                            g(Y)\ge1+N/10000
                           \qquad\text{for every }g(X)\ge0.
 \tag{QPL.16}
\]
The sphere case was proved above. An unramified elliptic
monodromy group would be abelian, so an elliptic base has
nontrivial branching. A nonscalar linear map fixes at most
the lines in spaces of dimensions \(d-1\) and one, giving
fixed fraction at most
\((q^{d-1}+q-2)/(q^d-1)\le2/3\).
Its index is at least \(N/6\), so \(g(Y)-1\ge N/12\).
For base genus at least two, Riemann--Hurwitz gives
\(g(Y)-1\ge N\). Both are stronger than (QPL.16).

\subsection*{3. All compatible field and graph operators}

The actual projective normalizer embeds in
\(\operatorname{Aut}(\operatorname{PSL}_d(q)^t)\):
perfectness removes a central discrepancy and Schur's lemma
makes the kernel scalar, exactly as in D-TD65.

The accepted degree barrier D-TD58 and the \(q\)-general
GLT degree/extension results give
\[
 a<q^{C_0d},\qquad 1\le j\le J_0<d/2,\qquad
 a>q^{d/12},\qquad \epsilon_q=q^{-B_0}.
 \tag{QPL.17}
\]
Here the inputs are Definition 2, Lemma 7.2, Theorem 7.3,
Corollary 8.6 and the complete dual-pair estimate (8.11) in
\cite{UNW:GLT2020}. They are stated for all prime powers.
The exact derivation in D-TD65 uses the field size, not
primality: \(j\ge d/2\) contradicts the degree barrier,
and \(a\ge q^{j(d-j)}/(q-1)\) bounds \(j\) otherwise.
The same exact error is
\(E\le4q^{2j^2-d}\le q^{-d/2}\).
The level-one vector/line induction argument also works
over every \(\mathbf F_q\). It follows that every compatible
inner-diagonal component operator has normalized trace
at most \(f_{\rm pt}+5q^{-d/2}\le f_{\rm pt}+\epsilon_q\).
No extension of its character to a field automorphism group
is assumed.

We give a separate count for every nontrivial field coset.
Let \(\sigma\ne1\) be an automorphism of \(\mathbf F_q\) and
\(\alpha(X)=M X^\sigma M^{-1}\).
Projectively fixed SL lifts satisfy
\(R(X)=cX\), where \(c\in\mathbf F_q^*\) and
\(R\) is \(\sigma\)-semilinear on
\(\operatorname{Mat}_d(\mathbf F_q)\).
Put \(\mathbf F_{q_0}=\operatorname{Fix}(\sigma)\).
For each \(c\), the fixed vectors of \(c^{-1}R\) which are
independent over \(\mathbf F_{q_0}\) are independent over
\(\mathbf F_q\).
Indeed a shortest \(\mathbf F_q\)-linear relation can be
normalized to have one coefficient one. Applying the
semilinear operator and subtracting leaves a shorter
relation, unless every coefficient is fixed by \(\sigma\).
The latter gives a forbidden \(\mathbf F_{q_0}\)-relation.

The solution space consequently has at most \(q_0^{d^2}\)
elements, whether or not \((c^{-1}R)^{\operatorname{ord}\sigma}=1\).
As \(q_0\le q^{1/2}\) and
\(|\operatorname{SL}_d(q)|>q^{d^2-1}/2\), we obtain
\[
             \frac{|C_{\operatorname{PSL}_d(q)}(\alpha)|}
                         {|\operatorname{PSL}_d(q)|}
                    \le2q^{-d^2/2+2}.
 \tag{QPL.18}
\]
This counts all allowed projective scalar discrepancies.

For a graph-field automorphism
\(\alpha(X)=M(X^\sigma)^{-T}M^{-1}\), including \(\sigma=1\),
the fixed lifts satisfy
\((X^\sigma)^TB_1X=cB_1\), with \(B_1=M^{-1}\).
For column \(j\), the first \(j-1\) equations from previous
columns are independent \(\mathbf F_q\)-linear constraints:
the transformed previous columns remain independent.
The complete column count therefore gives
\[
             \frac{|C_{\operatorname{PSL}_d(q)}(\alpha)|}
                         {|\operatorname{PSL}_d(q)|}
                    \le2q^{-d^2/2+d/2+2}.
 \tag{QPL.19}
\]
No symmetric, Hermitian or involutory form hypothesis is used.

Set \(R_0=\lfloor d/8\rfloor\).
As in the complete graph proof of D-TD65, a matrix of
geometric support less than \(R_0\) has a unique rational
largest eigenvalue and a rank factorization \(\lambda I+BC\).
There are at most \(q^{2dR_0+1}\) such projective elements.
Every nonempty fiber of
\(x\mapsto\alpha(x)x^{-1}\) has size \(|C_S(\alpha)|\).
Equations (QPL.18)--(QPL.19) bound the fraction of low-support
commutators by \(2q^{-d^2/8}\).

Use the accepted Larsen--Tiep all-element inequality (CFD.6)
with \(\sigma_0=10^{-29}\), the degree floor (QPL.17), and
the exact Schur average (CFD.10). Every compatible component
operator \(A\) inducing a non-inner-diagonal automorphism has
\[
 \frac{|\operatorname{tr}A|^2}{a^2}
       \le2q^{-d^2/8}+q^{-\sigma_0R_0/12}
       \le3q^{-2B_0-4}<q^{-2B_0}.
 \tag{QPL.20}
\]
The same integer inequalities checked in D-TD65 apply:
\(R_0\ge d/9\), \(d^2/8\ge2B_0+4\), and
\(d/(108\cdot10^{29})\ge2B_0+4\).
Thus all these outer operators have trace ratio below
\(\epsilon_q\), without extending a character through the
outer group.

\subsection*{4. An actual induced point action with linear frame stabilizers}

Let
\(\mathcal A=\operatorname{Aut}(\operatorname{PSL}_d(q))\)
and \(\mathcal B=\operatorname{PGL}_d(q)\).
Steinberg's Theorem 3.2 and its supplied normal-type proof
\cite[Sections 3--6]{PFL:Steinberg1960} give
\(\mathcal B\triangleleft\mathcal A\) and
\(c=[\mathcal A:\mathcal B]\le2f\).
The graph and field generators commute modulo \(\mathcal B\).
An exact semidirect splitting or equality \(c=2f\) is not needed.

Induce the projective-point \(\mathcal B\)-set to \(\mathcal A\).
It is a set of \(c\) frames, each containing \(N\) points.
Every element outside \(\mathcal B\) fixes no frame, because
the quotient acts regularly on the frames.
For an element of \(\mathcal B\), every frame has the same
point fixed fraction: field conjugation preserves rational
eigenspace dimensions, and inverse-transpose preserves the
point/dual-point fixed count.
This action is faithful. Its kernel lies in \(\mathcal B\)
from the frame action, and then vanishes from the faithful
point action in one frame.

Transport the action through the chosen actual component
identifications. The original group has an actual action on
a disjoint union \(\Omega\) of \(t\) such induced sets,
\[
          |\Omega|=ctN,\qquad c\le2f.
 \tag{QPL.21}
\]
This is the wreath action of the actual automorphism group
of the simple direct product, not a collection of quotient
maps to permuted individual components.

For each actual orbit, the stabilizer of a frame and component
has local image between PSL and PGL. Its associated connected
point cover has degree \(N\) over its actual intermediate base.
It is the same curve as the orbit's point-stabilizer quotient.
Equation (QPL.16) applies at that base's actual genus.
There is at least one orbit. Summing the original sphere
Riemann--Hurwitz equalities gives
\[
 \sum_i\frac{\operatorname{ind}_\Omega(y_i)}{ctN}
             \ge2+\frac2{10000ct}
             \ge2+\frac1{10000ft}.
 \tag{QPL.22}
\]

\subsection*{5. Whole partners and all original powers}

Let a layer operator have component permutation \(\pi\),
with \(s\) fixed components and \(I=t-\#\operatorname{cycles}(\pi)\).
The exact tensor trace identity bounds its normalized trace
by \(a^{-I}\) times the product of the normalized fixed-component
traces. Every return operator on a moved cycle remains present
and has normalized trace at most one.

If a fixed component is outer, its tiny trace pays the whole
product. Otherwise use the inner comparison and product
at most average. For \(s\ge1\),
\(a^I\ge1+2I\ge t-s+1\ge t/s\).
For \(s=0\), the trace is at most \(1/a<\epsilon_q\).
The frame fixed fractions just proved show, in all cases,
\[
             |\operatorname{tr}W(y)|/\dim W
                         \le f_\Omega(y)+\epsilon_q.
 \tag{QPL.23}
\]
Projectively trivial layer operators have ratio and fixed
fraction exactly one.

The accepted HD-FT20 fixed-component floor \(1/80\), the
HD-FT21 moved-cycle quarter floor, and tensor convolution
charge every nontrivial projective layer entry by \(e/80\).
Because \(e\ge a>480\), at most 160 original entries are active.
No floor on an omitted multiplicity partner is needed.

All trace comparisons use unitary representatives. Each finite
projective factor image has a finite preimage in its special
linear group; averaging a positive Hermitian form over this
finite lift supplies these representatives, also for the entire
partner \(U\). For the full original order \(L_i\) of \(T_i\),
its two unitary tensor factors have reciprocal scalar
\(L_i\)-th powers. Multiply them by reciprocal unit scalars
so that both \(L_i\)-th powers are one, leaving \(T_i\) unchanged.
The actual permutation period divides
\(L_i\). Every one of these original powers satisfies
(QPL.23), and the other entire factor has normalized trace
at most one. Fourier inversion at eigenvalue one therefore
gives actual moving rank at least
\(e(\operatorname{ind}_\Omega(y_i)/|\Omega|-\epsilon_q)\)
for an active entry. Retain the nonnegative actual rank for
all other entries. Equations (QPL.2) and (QPL.22) imply
\[
                \frac6e\ge\frac1{10000ft}-160\epsilon_q.
 \tag{QPL.24}
\]
But \(6/e<6\epsilon_q\).
Since \(f\le q\), \(q\ge2\), and \(B_0>1100000+24\),
\[
 166f\epsilon_q\le166q^{1-B_0}
       \le166\,2^{1-B_0}<2^{-1100000-14}
                         <\frac1{10000t}.
 \tag{QPL.25}
\]
This contradicts (QPL.24), proving the standalone theorem
for every proper prime power, without bounding \(f\).

For the exact stated minimal-source corollary, use the
already checked (PFL.17), namely \(N_*<T_0\), and the complete
D-TD54/D-TD12/D-TD16 interfaces exactly as in D-TD65.
All earlier and later factors remain in \(U\).
Later extracted cores are actual cores of quotient images;
they are not asserted to be subgroups of the original group.

The result excludes an unbounded linear natural-rank
residual with arbitrary whole partners and genuine field
normalizers. It does not classify unitary or other classical
low-level layers, bound original overgroups, or decide
source existence. Existing D-TD03--04 high-field coverage is
reused accepted information, not additional movement here.

